\documentclass[11pt]{article}
\usepackage[margin=1in]{geometry}
\usepackage{enumitem}
% =============================================================================
% Preambles/header.tex — shared LaTeX/Beamer preamble for this project
%
% Use from a Beamer lecture by prepending:
%     \input{header}
% and compiling with TEXINPUTS=../Preambles:$TEXINPUTS (the /compile-latex
% skill does this automatically).
%
% PALETTE CONTRACT. Color names defined here MUST match the names in
% Quarto/theme-template.scss so Beamer and Quarto renderings use the same
% palette. The scripts/check-palette-sync.sh script enforces this.
%
% To customize: edit the HEX values below AND the matching $variable in
% Quarto/theme-template.scss, then run scripts/check-palette-sync.sh.
% =============================================================================

% -----------------------------------------------------------------------------
% Engine detection + encoding
%
% Our /compile-latex skill uses XeLaTeX, where inputenc is unnecessary and
% can produce encoding warnings. Only load inputenc under pdfTeX.
% -----------------------------------------------------------------------------
\usepackage{iftex}
\ifPDFTeX
  \usepackage[utf8]{inputenc}
\fi

\usepackage{amsmath, amssymb, amsthm}
\usepackage{booktabs}
\usepackage{graphicx}

% -----------------------------------------------------------------------------
% PALETTE — mirrors Quarto/theme-template.scss variable names.
% Load xcolor and define colors BEFORE anything that references them
% (notably hyperref, hypersetup, and the Beamer color assignments).
% -----------------------------------------------------------------------------
\usepackage{xcolor}

\definecolor{primary-blue}{HTML}{012169}      % headings, accents
\definecolor{primary-gold}{HTML}{B9975B}      % emphasis, borders
\definecolor{highlight-yellow}{HTML}{F2A900}  % markers, alerts
\definecolor{light-bg}{HTML}{E8EDF5}          % title / section backgrounds
\definecolor{jet}{HTML}{1A1A1A}               % body text

% Semantic colors (used in diagrams and annotations)
\definecolor{positive}{HTML}{15803D}          % good / observed / identified
\definecolor{negative}{HTML}{B91C1C}          % bad / problematic / bias
\definecolor{neutral}{HTML}{525252}           % reference / context

% Highlight accents (match .hi-* classes in the SCSS)
\definecolor{hi-slate}{HTML}{314F4F}
\definecolor{hi-green}{HTML}{2E8B57}
\definecolor{hi-red}{HTML}{C41E3A}

% Semantic aliases so skill templates and snippets can write
% \textcolor{accent}{...} without hard-coding "primary-blue".
\colorlet{accent}{primary-blue}
\colorlet{accent2}{primary-gold}

% -----------------------------------------------------------------------------
% Hyperref — loaded AFTER xcolor + palette so link colors resolve.
% -----------------------------------------------------------------------------
\usepackage{hyperref}
\hypersetup{colorlinks=true, linkcolor=primary-blue, urlcolor=primary-blue,
            citecolor=primary-blue, pdfborder={0 0 0}}

% -----------------------------------------------------------------------------
% Course code — set once per deck (after \input{header}, before \begin{document})
% via \coursecode{CS401}. Shown in the footer on every slide so a deck's course
% is identifiable without opening the title page. Empty by default.
% -----------------------------------------------------------------------------
\makeatletter
\newcommand{\coursecode}[1]{\gdef\@coursecodetext{#1}}
\gdef\@coursecodetext{}
\makeatother

% -----------------------------------------------------------------------------
% Beamer theme assignments (only apply under Beamer).
%
% We detect Beamer via \@ifclassloaded{beamer}, which is the documented,
% non-internal way. This must be wrapped in \makeatletter/\makeatother
% because \@ifclassloaded contains an @-catcode character.
% -----------------------------------------------------------------------------
\makeatletter
\@ifclassloaded{beamer}{%
  \usefonttheme{professionalfonts}
  \setbeamercolor{structure}{fg=primary-blue}
  \setbeamercolor{frametitle}{fg=primary-blue}
  \setbeamercolor{title}{fg=primary-blue}
  \setbeamercolor{subtitle}{fg=primary-gold}
  \setbeamercolor{author}{fg=primary-blue}
  \setbeamercolor{institute}{fg=neutral}
  \setbeamercolor{date}{fg=primary-gold}
  \setbeamercolor{itemize item}{fg=primary-blue}
  \setbeamercolor{itemize subitem}{fg=primary-gold}
  \setbeamercolor{alerted text}{fg=primary-gold}
  \setbeamercolor{block title}{fg=white, bg=primary-blue}
  \setbeamercolor{block body}{bg=light-bg}
  % Semantic block variants: block=definition/concept (blue), exampleblock=
  % worked example (green/positive), alertblock=key takeaway or caution (gold).
  % Keep to <=2 colored boxes per slide across ALL three variants (INV-7).
  \setbeamercolor{block title example}{fg=white, bg=positive}
  \setbeamercolor{block body example}{bg=light-bg}
  \setbeamercolor{block title alerted}{fg=white, bg=primary-gold}
  \setbeamercolor{block body alerted}{bg=light-bg}

  % Minimal footer: course code (left) + page X / N (right), no navigation clutter
  \setbeamertemplate{navigation symbols}{}
  \setbeamertemplate{footline}{%
    \color{neutral}\footnotesize\hspace{0.7em}\@coursecodetext\hfill
    \insertframenumber/\inserttotalframenumber\hspace{0.7em}\vspace{0.3em}}
}{}
\makeatother

% -----------------------------------------------------------------------------
% TikZ setup — libraries the snippet gallery uses
% -----------------------------------------------------------------------------
\usepackage{tikz}
\usetikzlibrary{arrows.meta, positioning, calc, decorations.pathreplacing,
                fit, shapes.geometric, backgrounds}

% Shared TikZ styles — snippets and hand-written diagrams can pick these up
% via `\begin{tikzpicture}[every node/.style={...}]` or by naming specific
% styles (e.g., `\node[dag-node] (X) at (0,0) {X};`).
\tikzset{
  % Node styles for causal DAGs and similar
  dag-node/.style = {
    draw=primary-blue, fill=light-bg, rounded corners=2pt,
    minimum width=1.2cm, minimum height=0.9cm, align=center, font=\small
  },
  decision-node/.style = {
    draw=primary-gold, fill=white, diamond, aspect=1.8,
    minimum width=1.8cm, minimum height=1.2cm, align=center, font=\small,
    inner sep=1pt
  },
  % Edge styles — observed vs counterfactual
  observed-edge/.style = {-{Latex[length=2.5mm]}, thick, draw=jet},
  counterfactual-edge/.style = {-{Latex[length=2.5mm]}, thick, draw=neutral, dashed},
  confound-edge/.style = {-{Latex[length=2.5mm]}, thick, draw=negative, dashed},
  % Data-point styles for plots
  observed-dot/.style = {fill=primary-blue, draw=primary-blue, circle, inner sep=1.2pt},
  counterfactual-dot/.style = {fill=white, draw=neutral, circle, inner sep=1.2pt}
}

% -----------------------------------------------------------------------------
% Convenience macros
% -----------------------------------------------------------------------------
\newcommand{\muted}[1]{\textcolor{neutral}{#1}}
\newcommand{\key}[1]{\textcolor{primary-gold}{\textbf{#1}}}
\newcommand{\good}[1]{\textcolor{positive}{#1}}
\newcommand{\bad}[1]{\textcolor{negative}{#1}}

% Standout frame macro: full-bleed dark background for section transitions.
% Beamer-only (uses \begin{frame}/\setbeamercolor/\usebeamerfont) -- do not
% call from an article-class file (e.g. Notes/<CODE>/*-notes.tex). \muted,
% \key, \good, \bad, and \coursecode above are plain xcolor/gdef and are
% safe to use from either class; verified when Notes/article-class support
% was added.
% Expand: \transitionslide{Your transition title}
\newcommand{\transitionslide}[1]{%
  \begingroup
  \setbeamercolor{background canvas}{bg=primary-blue}
  \begin{frame}[plain]
    \centering
    \vfill
    {\usebeamerfont{title}\color{white}\Large #1\par}
    \vfill
  \end{frame}
  \endgroup
}

% Section divider frame (Beamer-only), an alternative to \transitionslide with a
% darker two-line style: near-black (jet) background, a small white label line
% above a larger gold title, and a thin gold rule along the bottom edge.
% Expand: \sectiondivider{Section 2}{Pointers}
\newcommand{\sectiondivider}[2]{%
  \begingroup
  \setbeamercolor{background canvas}{bg=jet}
  \begin{frame}[plain]
    \vspace*{3.0cm}
    {\color{white}\ttfamily\large #1\par}
    \vspace{0.5cm}
    {\color{primary-gold}\LARGE #2\par}
    \begin{tikzpicture}[remember picture, overlay]
      \draw[primary-gold, line width=2.5pt]
        ($(current page.south west)+(0,0.1)$) -- ($(current page.south east)+(0,0.1)$);
    \end{tikzpicture}
  \end{frame}
  \endgroup
}

% -----------------------------------------------------------------------------
% Channel watermark — "Concepts That Click" (YouTube).
%
% Article-class files (CompetitiveExam Q/A sets, Notes, Assignments) get a
% light diagonal draft watermark on every page plus a centered footer naming
% the channel. Beamer decks get a subtle TikZ background watermark instead
% (the footline already carries the course code). Centralized here so every
% MCQ PDF is branded without per-file edits.
% -----------------------------------------------------------------------------
\makeatletter
\@ifclassloaded{article}{%
  \usepackage{draftwatermark}
  \SetWatermarkText{Concepts That Click}
  \SetWatermarkScale{1}
  \SetWatermarkFontSize{2cm}
  \SetWatermarkLightness{0.86}
  \SetWatermarkAngle{45}
  \usepackage{fancyhdr}
  \pagestyle{fancy}
  \fancyhf{}
  \fancyfoot[C]{\footnotesize\textcolor{neutral}{Concepts That Click~|~YouTube: \href{https://www.youtube.com/@concepts.thatclick}{@concepts.thatclick}}}
  \renewcommand{\headrulewidth}{0pt}
  \renewcommand{\footrulewidth}{0pt}
  \fancypagestyle{plain}{%
    \fancyhf{}%
    \fancyfoot[C]{\footnotesize\textcolor{neutral}{Concepts That Click~|~YouTube: \href{https://www.youtube.com/@concepts.thatclick}{@concepts.thatclick}}}%
    \renewcommand{\headrulewidth}{0pt}%
    \renewcommand{\footrulewidth}{0pt}%
  }
}{}
\@ifclassloaded{beamer}{%
  \setbeamertemplate{background}{%
    \begin{tikzpicture}[remember picture, overlay]
      \node[rotate=45, opacity=0.055, align=center]
        at (current page.center)
        {\Huge\textcolor{neutral}{Concepts That Click}};
    \end{tikzpicture}%
  }
}{}
\makeatother 
\coursecode{JKSSB}
\renewcommand{\thesection}{36.\arabic{section}}
\newtheorem{example}{Example}[section]
\newenvironment{definitionbox}[1]{%
  \par\noindent\textbf{Definition (#1).}\ \itshape
}{\par\vspace{0.3em}}
\title{Web Browsing, Search, Upload and Download --- Detailed Lecture Notes}
\author{JKSSB Computer Awareness}
\date{\today}

\begin{document}
\maketitle

\section{The browser, tabs, bookmarks, and history}
A \textbf{web browser} is the software used to open and view web pages on the
internet. Common examples are Google Chrome, Mozilla Firefox, and Microsoft
Edge. The browser sends a request for a page and then displays the returned
text, images, and links. The address typed to reach a page is the \textbf{URL}
--- Uniform Resource Locator --- and it is typed in the \textbf{address bar}
at the top of the browser window.

\begin{definitionbox}{Web browser}
Software that requests web pages from servers and displays them to the user.
\end{definitionbox}

\begin{center}
\begin{tabular}{p{0.24\textwidth}p{0.66\textwidth}}
\toprule
\textbf{Term} & \textbf{Meaning in this lesson} \\
\midrule
Web browser & Software (Chrome, Firefox, Edge) that opens and displays web pages. \\
Tab & One open page inside the browser window; many tabs can stay open together. \\
Bookmark & A saved shortcut to a page for quick return without retyping the address. \\
Browsing history & The date-wise list of pages visited earlier, kept by the browser. \\
URL & Uniform Resource Locator: the address of a page, typed in the address bar. \\
Search engine & A website (Google, Bing) that finds pages matching typed words. \\
\bottomrule
\end{tabular}
\end{center}

Each open page lives in its own \textbf{tab}, so several pages can stay open
in one window. A new tab is opened with \textbf{Ctrl+T}. A \textbf{bookmark}
is a saved shortcut to a page: instead of typing the address again, the saved
bookmark is clicked. It is created with \textbf{Ctrl+D}. \textbf{Browsing
history} is the list of pages visited earlier, stored by date, and it is
opened with \textbf{Ctrl+H}. The distinction is simple: history answers
``where was I'', while a bookmark answers ``where do I want to return''.

\begin{example}
While preparing a report on the office desktop, the clerk keeps the office
portal open in one tab and the search results in a second tab, and bookmarks
the portal page so it can be opened directly the next morning. The history
list that evening shows both pages in order.
\end{example}

\section{How a search engine works}
A \textbf{search engine} is a website, such as Google or Bing, that finds
pages matching the words typed by the user. It works in three ordered stages:
crawling, indexing, and ranking.

\begin{definitionbox}{Crawling}
The collection step: automatic programs called crawlers or spiders visit
pages and follow their links to discover new and changed pages.
\end{definitionbox}

\begin{definitionbox}{Indexing}
The storage step: the crawled pages are kept in a large index, like the
index at the back of a book, so matches can be found quickly.
\end{definitionbox}

\begin{definitionbox}{Ranking}
The ordering step: when words are typed, matching pages are picked from the
index and ordered by relevance on the search results page.
\end{definitionbox}

Crawling comes first because without it the engine would never learn that a
new page exists. Indexing comes second: the collected pages are stored in the
index before any search is typed. Ranking comes last: only after the user
types words does the engine select matches from that ready index and order
them. The anchor to remember is the order itself: \textbf{crawl, then index,
then rank}.

\section{Sponsored results and organic results}
The matches chosen by relevance are called \textbf{organic results}. In
addition, the results page shows \textbf{sponsored results}: paid placements
that appear because an advertiser paid for them, not because they ranked
highest on relevance. A sponsored result always carries a visible
\textbf{Ad} or \textbf{Sponsored} label. The practical consequence is direct:
being listed first on the page does not always mean being the most relevant
match, since the top position may be a paid placement.

\begin{example}
A search for a printer shows two blocks: the first carries a small ``Ad''
label and links to a seller who paid for the position, while the entries
below it are the organic results ordered by relevance. The topmost entry is
therefore not necessarily the best match.
\end{example}

\section{Cookies, cache, and history}
These three stores are the most confused triple in this lesson. Each stores a
different thing, and each is stored by a different party.

\begin{definitionbox}{Cookie}
A small file a website stores on the computer to remember that visitor:
login state, cart items, and preferences such as language.
\end{definitionbox}

\begin{definitionbox}{Browser cache}
Copies of page parts (images, styles, scripts) kept by the browser on the
computer so repeat visits load faster.
\end{definitionbox}

\begin{definitionbox}{Browsing history}
The list of visited page addresses, kept by the browser by date.
\end{definitionbox}

\begin{center}
\begin{tabular}{p{0.18\textwidth}p{0.36\textwidth}p{0.36\textwidth}}
\toprule
\textbf{Item} & \textbf{What it stores} & \textbf{Who stores it} \\
\midrule
Cookie & Visitor data: login, cart, choices & The website \\
Cache & Copies of page files: images, styles & The browser \\
History & List of visited page addresses & The browser \\
\bottomrule
\end{tabular}
\end{center}

Cookies are set by websites and read back on the next visit, which is why a
user stays logged in. The cache is managed by the browser: on the next visit
the saved copies load from the disk instead of being fetched again, so the
page opens faster. Clearing them has different effects. Deleting
\textbf{cookies} signs the user out of sites and clears saved preferences.
Clearing the \textbf{cache} does not sign the user out; it only makes the
next visit slightly slower. Clearing \textbf{history} removes the record of
visited addresses but changes neither logins nor saved page copies.

\section{Downloading and uploading}
\textbf{Downloading} means copying a file from the internet to the user's
computer: saving a form PDF from the office portal, or saving a photo from a
web page. \textbf{Uploading} means sending a file from the user's computer to
the internet: attaching a document to an email, or submitting a scanned form
on a portal. The direction is always judged from the user's machine: download
moves data \textbf{down} to the machine, upload moves data \textbf{up} from
the machine to a server.

\begin{definitionbox}{Downloading}
Copying a file from the internet (a server) to the user's computer.
\end{definitionbox}

\begin{definitionbox}{Uploading}
Sending a file from the user's computer to the internet (a server).
\end{definitionbox}

\begin{center}
\begin{tabular}{p{0.22\textwidth}p{0.34\textwidth}p{0.34\textwidth}}
\toprule
\textbf{Point} & \textbf{Download} & \textbf{Upload} \\
\midrule
Direction & Server to computer & Computer to server \\
Typical use & Opening pages, saving files & Mail attachments, form submission \\
Typical speed & Faster & Slower \\
\bottomrule
\end{tabular}
\end{center}

Home and office connections usually provide a faster download speed than
upload speed, because receiving pages and files is the more common need. A
downloaded file, once saved, opens without needing the internet; an upload
needs the connection at the moment of sending.

\section{File conversion}
\textbf{File conversion} means changing a file from one format to another
while keeping the same content. The common office need is converting a Word
document (\textbf{.docx}) to \textbf{PDF} before submitting it on a portal.
Conversion changes the \textbf{format}, not the transfer direction: a
converted file must still be uploaded or downloaded separately. The
\textbf{extension} (such as .pdf, .docx, .jpg) is the name tag that says
which format the file is stored in, while the \textbf{format} (such as PDF,
PNG) is the storage method itself. The two words must not be conflated: the
extension names the format; it is not the format.

\begin{example}
The clerk finishes the report in Word, converts the .docx file to PDF, and
then uploads the PDF on the portal. The conversion changed the format; the
upload changed the location. If the portal instead needed a photograph, a
.JPG image would be converted to the needed picture format the same way.
\end{example}

\section{Exam retrieval and common confusions}
Five distinctions prevent most errors on this topic.
\begin{enumerate}
  \item The search order is \textbf{crawl}, then \textbf{index}, then
    \textbf{rank}. The index is built before the search is typed.
  \item \textbf{Sponsored} means a paid placement with an Ad label;
    \textbf{organic} means ranked by relevance. The top of the page is not
    always the best match.
  \item \textbf{Cookie} (website; visitor data) vs \textbf{cache} (browser;
    page copies) vs \textbf{history} (browser; visited list). Do not swap
    the three.
  \item \textbf{Download} means internet to computer; \textbf{upload} means
    computer to internet. Judge the direction from the user's machine.
  \item \textbf{Conversion} changes the format; \textbf{transfer} changes
    the location. The extension (.png) names the format (PNG); it is not
    the format itself.
\end{enumerate}

\subsection*{Exercises}
\begin{enumerate}[label=\arabic*.]
  \item List the three stages of a search engine in order and say what each stage does.
  \item Distinguish a sponsored result from an organic result.
  \item Distinguish a cookie from the browser cache and from browsing history.
  \item Explain the difference between downloading and uploading, giving one office example of each.
  \item What does file conversion change, and how is the extension related to the format?
\end{enumerate}

\subsection*{Solutions}
\begingroup\small
\begin{enumerate}[label=\arabic*.]
  \item Crawling (crawlers visit pages and follow links to collect them),
    then indexing (pages stored in a large index), then ranking (matches
    picked from the index and ordered by relevance after words are typed).
  \item A sponsored result is a paid placement carrying an Ad or Sponsored
    label; an organic result is a normal match chosen and ordered by
    relevance.
  \item A cookie is visitor data (login, cart, choices) stored by the
    website; the cache is copies of page files (images, styles) stored by
    the browser for faster repeat visits; history is the browser's
    date-wise list of visited addresses.
  \item Downloading copies a file from the internet to the computer, such as
    saving a form PDF from the portal; uploading sends a file from the
    computer to the internet, such as attaching a document to an email.
  \item Conversion changes the format while keeping the content (for example
    .docx to PDF). The extension (.pdf, .docx) is the name tag for the
    format; the format is the storage method itself.
\end{enumerate}
\endgroup
\end{document}
